\section{A finite-energy estimate at the moving marked point}
\label{eng:section}

In this section the approximating domains have a common physical diameter
bound.  We work with any of their subsequential interface laws supplied by
Proposition~\ref{geo:compactness}.  Constants may depend on this bound,
the fixed mass, the common rhombic angle bound, and the resulting crossing
constants.  They do not depend on a test point, a stopping time, or the
particular subsequential law.  The natural filtration of the continuous
driver, with its usual augmentation, is denoted by $(\cF_t)$.

Only geometric crossing estimates and deterministic Loewner identities
are used below.  In particular, we do not yet assume that the driver is a
semimartingale.  The distinction is consequential: the estimate proved
here will justify the stochastic reconstruction of the driver in the next
section.

Recall the initial chart $\phi:\Omega\to\Hh$, the half-plane Loewner maps
$g_t$, and the translated inverse chart
\[
 \psi_t(v)=\phi^{-1}\bigl(g_t^{-1}(W_t+v)\bigr),
 \qquad \mu_t(v)=m|\psi_t'(v)|.
\]
All integrals in this section are nonnegative integrals and are allowed
initially to take the value $+\infty$.

\begin{proposition}[Moving-mark energy]\label{eng:energy}
For
\begin{equation}\label{eng:J-definition}
 J_t=\int_{\Hh}\frac{\mu_t(v)}{|v|^2}\,\dd A(v),
\end{equation}
there is a constant $K<\infty$ such that
\begin{equation}\label{eng:energy-bound}
 \E\int_0^\infty J_t^2\,\dd t
 \le K m^2\area(\Omega)<\infty.
\end{equation}
Consequently $J_t$ is finite for almost every capacity time, almost surely.
\end{proposition}

We prove the proposition by following one physical point in a clock
adapted to its distance from the driver.  Keeping a small positive power
of its angular coordinate is necessary when summing over physical area.

\subsection{The point clock and the two geometric inputs}

Fix $w\in\Omega$.  Its lifetime $T_w$ is the supremum of the capacity times
for which $w$ belongs to the active domain $O_t$.  Before that lifetime put
\begin{equation}\label{eng:point-variables}
 Z_t=g_t(\phi(w))-W_t,\qquad
 r_t=|Z_t|,\qquad Y_t=\Ima Z_t,\qquad X_t=\frac{Y_t}{r_t}\in(0,1].
\end{equation}
The letter $r_t$ in this section denotes a radial coordinate and is
unrelated to the elliptic nome $r_\delta$.
Define
\[
 \sigma(t)=\int_0^t\frac{\dd u}{r_u^2},\qquad 0\le t<T_w,
 \qquad
 \sigma_* =\lim_{t\uparrow T_w}\sigma(t),
\]
where $T_w=\infty$ is permitted.  Let $t(s)$ be the inverse clock on
$0\le s<\sigma_*$.  A subscript $s$ on $X$, $Y$, or $r$ will mean evaluation
at $t(s)$ whenever we explicitly work in this clock.  Set $X_s=0$ for
$s\ge\sigma_*$.  This killing convention applies also if $T_w=\infty$
but the total clock $\sigma_*$ is finite.  Ratios involving $r_s$ below are
used only for $s<\sigma_*$, and their entire integrands are set to zero
afterward.

\begin{lemma}[Deterministic clock identities]\label{eng:clock}
Before the point's lifetime,
\begin{equation}\label{eng:clock-identities}
 Y_s=Y_0e^{-2s},\qquad
 \dd\log\crad_{O_{t(s)}}(w)=-4X_s^2\,\dd s,
 \qquad \dd t=r_s^2\,\dd s.
\end{equation}
The same identities hold after any live stopping time, with elapsed
capacity and clock times and the corresponding starting values.
\end{lemma}

\begin{proof}
The imaginary part of the Loewner equation gives
$\partial_tY_t=-2Y_t/r_t^2$.  This differentiates $g_t(\phi(w))$, not
$W_t$: the driver is real, so its subtraction does not affect the
imaginary part.  If $H_t=g_t\circ\phi$, then
\[
 \crad_{O_t}(w)=\frac{2Y_t}{|H_t'(w)|},\qquad
 \partial_t\log|H_t'(w)|=-2\Rea\frac1{Z_t^2}.
\]
It follows that
\[
 \partial_t\log\crad_{O_t}(w)
 =-\frac2{r_t^2}+2\Rea\frac1{Z_t^2}
 =-\frac{4Y_t^2}{r_t^4}.
\]
The clock has a strictly positive continuous rate on every compact live
interval.  Changing variables proves all the assertions.
\end{proof}

We spell out the precise restarted estimates inherited from the discrete
crossing bounds.  The first is a boundary estimate in a current conformal
chart; the second is a bulk estimate in physical coordinates.  The latter
uses the FK rectangle bounds as well as the conditional geometry.  It is
not an application of the boundary form of G3 to an annulus wholly inside
the domain.

\begin{lemma}[Restarted geometric estimates]\label{eng:restarts}
There are $a>0$, $K<\infty$, and $d_*>0$ with the following properties.
They hold conditionally on the past at every stopping time $\tau$ on the
event $w\in O_\tau$.
\begin{enumerate}
\item Translate and dilate the current half-plane chart so that its tip
is $0$, its target is $\infty$, and $|Z_\tau|=1$.  If
$X_\tau\le d_*$ and $x=\Rea(Z_\tau/|Z_\tau|)$, then, for
$0<d<d_*$, the probability that the restarted trace visits $B(x,d)$ is at
most $Kd^a$.
\item With $R=\crad_{O_\tau}(w)$ and $0<\varepsilon<1$,
\begin{equation}\label{eng:bulk-crad}
 \Prob\left\{
   \inf_{\tau\le t<T_w}\crad_{O_t}(w)<\varepsilon R
      \,\middle|\,\cF_\tau\right\}
 \le K\varepsilon^a.
\end{equation}
The infimum in this formula is strictly before the lifetime.  Conformal
radius is not assigned the value zero upon swallowing.
\end{enumerate}
These estimates remain valid for the inverse-clock stopping times and
the first angular-level stopping times used below.
\end{lemma}

\begin{proof}
For the first assertion, decrease $d_*$ so that $|x|$ stays, for example,
between $3/4$ and $1$.  A half-annulus centered at $x$, with a fixed outer
radius less than $1/4$ and inner radius $d$, misses the tip $0$ and the
target $\infty$ and is unforced.  A visit to its inner ball entails an
unforced crossing.  Conditional G3 and its uniform conformal transport
therefore give the asserted bound.  A subsequent swallowing of the ball
without visiting it is not a crossing and causes no difficulty for this
statement.

For the second assertion, Koebe's conformal-radius comparison gives
\[
 \operatorname{dist}(w,\partial O_\tau)\ge R/4.
\]
If $w$ is still alive and its conformal radius drops below
$\varepsilon R$, its distance to the current boundary is less than
$\varepsilon R$.  For sufficiently small $\varepsilon$, this new nearby
boundary cannot be a piece of the old boundary, which is at distance at
least $R/4$.  The future interface must therefore approach within
$\varepsilon R$ of $w$.  Such an approach supplies an actual FK arm
through the annulus with inner radius comparable to $\varepsilon R$ and
outer radius, say, $R/8$.  This annulus is initially in the protected bulk.
In each of a sequence of disjoint buffered shells, rectangle crossings
and FK positive association give a blocking circuit with a probability
bounded below, uniformly under every conditioning outside that shell.
Their product gives $K\varepsilon^a$.  The physical outer radius is
uniformly bounded by the assumed diameter bound.  Larger values of
$\varepsilon$ are absorbed by increasing $K$.

Here these statements are inherited by the subsequential law through
the localized conditional convergence of Corollary~\ref{geo:restart}.
We detail the role of the stopping times, since no stochastic
differential equation for $W$ is available.  First restrict to a finite
capacity interval on which the test point has positive chart height,
its relevant coordinates stay bounded, and its physical access and
conformal radius are protected.  Driver convergence and the Loewner ODE
then imply convergence of its coordinates, derivatives, and point
clocks.  The inverse clock is continuous on such an interval because
its rate is positive and continuous.  In the discrete conditional
future, finish the current tile when necessary and use the uniform
crossing bounds; the tile error disappears under the same protections.
One passes these bounds by multiplying them by bounded continuous
functions of the stopped driver history, slightly enlarging visiting
balls and relaxing conformal-radius thresholds before taking the limit.
The conditional inequalities consequently hold in the natural driver
filtration.

For a bounded live stopping time, approximation from above by stopping
times on a deterministic time grid gives the same conclusion.  Current
charts and coordinates are continuous under the protections; relaxed
thresholds remove a possible equality at the limiting time.  First
angular-level times can equivalently be approximated on clock grids,
with a slightly relaxed level, and then obtained by continuity.
Finally exhaust the finite horizons and the protections.  For the
physical annulus estimate, take the mesh to zero at each fixed positive
protected radius before removing that protection.  Thus the annuli used
in the limiting estimate are always above the lattice cutoff during its
proof.  This establishes the two estimates, including all their live
clock restarts, without assuming a continuum domain Markov property or
a semimartingale driver.  Only the natural past is conditioned on; no
enlargement revealing future information is used.
\end{proof}

\subsection{A deterministic estimate on boundary avoidance}

\begin{lemma}[Reflected-ball estimate]\label{eng:reflection}
Consider a chordal Loewner trace in $\Hh$ with continuous driver, starting
from $0$.  Let its maps be $l_t$, so that $l_0$ is the identity.  Start a
point at $z_0=x+\ii X_0$ with $|z_0|=1$.  There are absolute constants
$c_0>0$ and $K<\infty$ such that, if $0<d<1/8$, $X_0\le c_0d$, and the trace avoids
$B(x,d)$, then up to the point's lifetime
\begin{equation}\label{eng:avoidance-deterministic}
 \sup_s\frac{X_s}{X_0}\le\frac Kd,
 \qquad
 \int_0^{\sigma_*}\frac{X_s}{X_0}\,\dd s\le\frac Kd.
\end{equation}
Consequently
\begin{equation}\label{eng:avoidance-square}
 \int_0^{\sigma_*}\left(\frac{X_s}{X_0}\right)^2\dd s
 \le\frac K{d^2}.
\end{equation}
\end{lemma}

\begin{proof}
Shrink the ball radius by a fixed factor whenever needed below; this
only changes the constants.  Because the trace misses the original
ball, its connected upper half-ball is disjoint from the trace at every
time.  All points of the smaller upper half-ball are in the same
component of the trace complement.  They are therefore disconnected
from infinity simultaneously, if disconnection occurs.  The real
boundary germ at $x$ will have this same swallowing time, as we verify
below.  This assertion uses avoidance of a whole ball, not merely
avoidance of the point $x$.

Before this common lifetime the map extends by Schwarz reflection to a
univalent map on a smaller full ball centered at $x$.  Its derivative
$l_t'(x)$ is positive.  The driver is outside this reflected image:
if $W_t=l_t(\xi)$ for a point in the ball, symmetry and univalence
force $\xi$ to be real.  The local inverse would then give
\[
 \gamma(t)=\lim_{y\downarrow0}l_t^{-1}(W_t+\ii y)=\xi,
\]
contrary to trace avoidance.

Fix any $t_0$ before the common bulk lifetime. Since
$l_{t_0}(x)\ne W_{t_0}$, the complex Loewner ODE has a flow
$\Phi_{t,t_0}$, holomorphic near $l_{t_0}(x)$, for $t$ in a time
neighborhood of $t_0$ within this lifetime (one-sided at $0$).
On a sufficiently small fixed disc $D$ about $x$, the composition
$\Phi_{t,t_0}\circ l_{t_0}$ is defined and holomorphic.
Forward and backward ODE uniqueness identify it with $l_t$ on
$D\cap\Hh$; the identity theorem extends this equality to $D$.
Hence $t\mapsto l_t(x)$ satisfies
$\partial_t l_t(x)=2/(l_t(x)-W_t)$ before the common bulk lifetime.

Thus the real Loewner solution at $x$
cannot meet the driver before the common bulk lifetime.  Conversely,
if that real solution survived a finite bulk disconnection time, its
distance from the continuous driver would be positive on a slightly
longer compact time interval.  Continuous dependence for the complex
Loewner ODE would then make all sufficiently close points $x+\ii y$
survive that longer interval.  This contradicts the simultaneous
disconnection of all points in the half-ball.  The real and bulk
lifetimes therefore agree; the argument also covers an infinite
lifetime by exhaustion.

Koebe's theorem and
interior distortion yield, with
$r_x(t)=|l_t(x)-W_t|$,
\begin{equation}\label{eng:reflection-distortion}
 r_x(t)\ge c d\,l_t'(x),\qquad
 \left|\frac{l_t''(x)}{l_t'(x)}\right|\le\frac Kd,
 \qquad
 Y_t\asymp X_0l_t'(x),\qquad r_t\asymp r_x(t).
\end{equation}
For the last comparison, distortion first gives
$|l_t(z_0)-l_t(x)|\le KX_0l_t'(x)$; choosing $c_0$ sufficiently small
and using the first inequality makes this a small fraction of $r_x(t)$.

Differentiate the deterministic Loewner equation in its spatial
variable.  On the real reflected interval it gives
\begin{equation}\label{eng:logarithmic-derivative}
 \partial_t\left(\frac{l_t''}{l_t'}\right)(x)
   =\frac{4l_t'(x)}{(l_t(x)-W_t)^3}.
\end{equation}
The real denominator cannot change sign before swallowing.  Hence the
right-hand side has a fixed sign.  The primitive starts from zero and
is bounded in absolute value by $K/d$ through
\eqref{eng:reflection-distortion}.  It follows that
\begin{equation}\label{eng:real-integral}
 \int\frac{l_t'(x)}{r_x(t)^3}\,\dd t\le\frac Kd
\end{equation}
up to any strictly earlier time and then up to the lifetime by monotone
convergence.  This argument differentiates a map with respect to space
and capacity; it never differentiates the continuous driver.

The distortion inequalities give
$X_t/X_0\asymp l_t'(x)/r_x(t)\le K/d$.  Since
$\dd s=\dd t/r_t^2$ and $r_t\asymp r_x(t)$,
\[
 \int\frac{X_s}{X_0}\,\dd s
 \le K\int\frac{l_t'(x)}{r_x(t)^3}\,\dd t\le\frac Kd.
\]
Taking the product of this bound and the supremum bound proves
\eqref{eng:avoidance-square}.  Disconnection introduces no additional
term: the integrals stop at the common lifetime, and the killed angular
process is zero beyond its clock lifetime.
\end{proof}

\subsection{Uniform angular controls and exponential contraction}

\begin{lemma}[Supremum tails and occupation tightness]\label{eng:controls}
At every live restart, use elapsed point clock and write $X_0$ for its
starting angular coordinate.  There are $a>0$ and $K<\infty$ such that
\begin{equation}\label{eng:supremum-tail}
 \Prob\left\{\sup_{s\ge0}\frac{X_s}{X_0}>u
       \,\middle|\,\text{past}\right\}\le Ku^{-a},\qquad u\ge1.
\end{equation}
Moreover the conditional laws of
\begin{equation}\label{eng:occupation}
 I=\int_0^\infty\left(\frac{X_s}{X_0}\right)^2\dd s
\end{equation}
are uniformly tight, over all live restarts and all starting angles.
\end{lemma}

\begin{proof}
Translate and dilate the current chart to make the starting radial
coordinate equal to one.  The corresponding capacity scaling leaves
$\dd t/r_t^2$, and hence the point clock, unchanged.
Time zero and the notation $O_0$ in this proof refer to the chosen
restart.

For a large $u$ take $d=A/u$, with a fixed $A$ larger than both the
constant in \eqref{eng:avoidance-deterministic} and $1/c_0$.  If
$X_0\le c_0d$, the reflected-ball bound makes the supremum at most $u$
unless the trace visits $B(x,d)$.  Lemma~\ref{eng:restarts} bounds that
exception by $Kd^a$.  If $X_0>c_0d$, the deterministic inequality $X\le1$
already gives $X/X_0<u$.  Increasing $K$ handles the bounded range of
$u$ for which $d$ is not yet sufficiently small.  This proves
\eqref{eng:supremum-tail}.

For tightness, fix a sufficiently small $d>0$.  If $X_0\le c_0d$, then
$I\le K/d^2$ outside an event of conditional probability at most $Kd^a$,
by Lemmas~\ref{eng:restarts} and~\ref{eng:reflection}.  If
$X_0>c_0d$, the conformal-radius identity instead gives
\[
 I=\frac1{4X_0^2}
    \log\frac{\crad_{O_0}(w)}
                  {\inf_{t<T_w}\crad_{O_t}(w)}.
\]
This is an equality of extended nonnegative quantities, with the
infimum taken on the live interval.  By \eqref{eng:bulk-crad},
\[
 \Prob\{I>L\mid\text{past}\}
 \le K\exp(-4a c_0^2d^2L)
 \qquad\text{if }X_0>c_0d.
\]
For $L\ge K/d^2$ the two cases therefore give a uniform bound by
$Kd^a+K\exp(-4a c_0^2d^2L)$.  Given an error probability, first choose
$d$ small and then $L$ large.  This proves the asserted tightness.
\end{proof}

\begin{lemma}[Small-power contraction]\label{eng:decay}
There exist $p\in(0,1/4)$, $b>0$, and $K<\infty$ such that, from every
live restart,
\begin{equation}\label{eng:moment-decay}
 \E[X_s^p\mid\text{past}]\le K X_0^p e^{-bs},\qquad s\ge0.
\end{equation}
The same choice of $p$ has a uniformly bounded conditional $2p$-moment
of $\sup_{s\ge0}X_s/X_0$.
\end{lemma}

\begin{proof}
Choose $p<1/4$ so small that $4p<a$, with $a$ from
\eqref{eng:supremum-tail}.  Integrating that tail gives
\begin{equation}\label{eng:supremum-moment}
 \E\left[\left(\sup_{s\ge0}\frac{X_s}{X_0}\right)^{2p}
          \,\middle|\,\text{past}\right]\le K_p.
\end{equation}
We next find a deterministic clock time $H$ at which the conditional
$p$-moment contracts strictly.

Fix a small error tolerance $\alpha>0$, to be chosen below.  Choose
$\rho\in(0,1/2)$ so small that $K(2\rho)^a\le\alpha$.  By occupation
tightness choose $L$ with
$\Prob\{I>L\mid\text{past}\}\le\alpha$, uniformly over the restart.
Take $H>L/\rho^2$.  If the angular process is alive throughout $[0,H]$
and never reaches $\rho X_0$, then $I\ge\rho^2H>L$.  Thus with
conditional probability at least $1-\alpha$ it either dies or reaches
$\rho X_0$ by $H$.  At the first live hitting time, the restarted
supremum estimate gives
\[
 \Prob\left\{\sup_{\text{after hit}}X_s>X_0/2
        \,\middle|\,\text{past at hit}\right\}
 \le K(2\rho)^a\le\alpha.
\]
The killed process never rebounds.  Consequently
\begin{equation}\label{eng:half-probability}
 \Prob\{X_H/X_0>1/2\mid\text{past}\}\le2\alpha.
\end{equation}
Using \eqref{eng:supremum-moment} and conditional Cauchy--Schwarz on the
exceptional event,
\[
 \E[(X_H/X_0)^p\mid\text{past}]
 \le 2^{-p}+K_p^{1/2}(2\alpha)^{1/2}.
\]
Choose $\alpha$ sufficiently small to make the right-hand side a
constant $\lambda<1$.  All choices were uniform in the starting angle
and past.

Apply this contraction successively at clock times $H,2H,\ldots$ while
alive; after death the conditional contribution is zero.  This gives
$\E[X_{nH}^p\mid\text{initial past}]\le\lambda^nX_0^p$.  For
$s=nH+u$, $0\le u<H$, the restarted supremum moment bounds the last
partial interval by a fixed factor.  With
$b=-H^{-1}\log\lambda>0$, enlarging $K$ proves
\eqref{eng:moment-decay}.  Only this $p$-moment is asserted to contract;
the higher moment needed later comes from
\eqref{eng:supremum-moment}.
\end{proof}

\subsection{The two capacity-to-clock ratios}

\begin{lemma}[Weighted clock integral]\label{eng:weighted}
For the $p$ in Lemma~\ref{eng:decay}, there are $\epsilon>0$ and
$K<\infty$ such that, uniformly over every initial live point,
\begin{equation}\label{eng:weighted-clock}
 \E\int_0^{\sigma_*}X_s^{2p}
 \left[
   \left(\frac{\sqrt{t(s)}}{r_s}\right)^\epsilon
  +\left(\frac{r_s}{\sqrt{t(s)}}\right)^\epsilon
 \right]\dd s\le K.
\end{equation}
Here $t(s)$ is capacity elapsed from the chosen starting instant.  The
value at $s=0$ is immaterial.
\end{lemma}

\begin{proof}
On a surviving clock interval, Lemma~\ref{eng:clock} gives
\begin{equation}\label{eng:time-formula}
 r_s=\frac{Y_0e^{-2s}}{X_s},\qquad
 t(s)=Y_0^2\int_0^s\frac{e^{-4v}}{X_v^2}\,\dd v.
\end{equation}
Since $X_v\le1$, the second integral is at least a fixed constant times
$\min(1,s)$.  Therefore
\begin{equation}\label{eng:inverse-ratio}
 \frac{r_s}{\sqrt{t(s)}}
 \le\frac{K}{X_s\sqrt{\min(1,s)}}.
\end{equation}
For $0<\epsilon\le p$, the inequality $X_s^{2p-\epsilon}\le X_s^p$
and \eqref{eng:moment-decay} imply
\begin{equation}\label{eng:inverse-ratio-mean}
 \E\left[X_s^{2p}
       \left(\frac{r_s}{\sqrt{t(s)}}\right)^\epsilon;
          s<\sigma_*\right]
 \le K\min(1,s)^{-\epsilon/2}X_0^p e^{-bs}.
\end{equation}
Its integral is uniformly finite when $\epsilon<2$.

For the other ratio put
$q_s=\min_{0\le v\le s}X_v$ on $\{s<\sigma_*\}$.  This minimum is
strictly positive on each surviving compact clock interval.
Equation~\eqref{eng:time-formula} gives
\begin{equation}\label{eng:forward-ratio}
 \frac{\sqrt{t(s)}}{r_s}
 \le\sqrt{s}\,e^{2s}\frac{X_s}{q_s}.
\end{equation}
Partition the surviving event into the disjoint bins
\[
 B_k(s)=\{s<\sigma_*,\quad
          2^{-k-1}X_0<q_s\le2^{-k}X_0\},\qquad k=0,1,\ldots.
\]
For $k\ge1$, on $B_k(s)$ the process has reached $2^{-k}X_0$ while
alive by clock time $s$.  Restarting the $2p$-moment supremum estimate
there gives
\begin{equation}\label{eng:bin-hit}
 \E[X_s^{2p};B_k(s)]\le K X_0^{2p}2^{-2pk}.
\end{equation}
For $k=0$ the same bound follows by restarting at time zero.  Independently,
$X_s^{2p}\le X_s^p$ and the small-power contraction give
\begin{equation}\label{eng:bin-decay}
 \E[X_s^{2p};B_k(s)]\le K X_0^p e^{-bs}.
\end{equation}
Taking the geometric mean of these two upper bounds yields
\begin{equation}\label{eng:bin-interpolation}
 \E[X_s^{2p};B_k(s)]
 \le K X_0^{3p/2}2^{-pk}e^{-bs/2}.
\end{equation}
No exponential contraction of the $2p$-moment has been used.

On $B_k(s)$ we have
$q_s^{-\epsilon}\le2^{(k+1)\epsilon}X_0^{-\epsilon}$.
Combining this with \eqref{eng:forward-ratio},
$X_s^{2p+\epsilon}\le X_s^{2p}$, and
\eqref{eng:bin-interpolation}, we obtain the full bin sum
\begin{align}
 &\E\left[X_s^{2p}
       \left(\frac{\sqrt{t(s)}}{r_s}\right)^\epsilon;
                  s<\sigma_*\right] \notag\\
 &\qquad\le K s^{\epsilon/2}e^{2\epsilon s}X_0^{-\epsilon}
       \sum_{k\ge0}2^{k\epsilon}\E[X_s^{2p};B_k(s)] \notag\\
 &\qquad\le
 K s^{\epsilon/2}e^{-(b/2-2\epsilon)s}
       X_0^{3p/2-\epsilon}
       \sum_{k\ge0}2^{-(p-\epsilon)k}.
 \label{eng:bin-sum}
\end{align}
Choose once and for all
\begin{equation}\label{eng:epsilon-choice}
 0<\epsilon<\min\{p,b/4,2\}.
\end{equation}
The geometric series converges, $X_0^{3p/2-\epsilon}\le1$, and the
remaining expression is integrable at both zero and infinity.
Together with \eqref{eng:inverse-ratio-mean}, this proves
\eqref{eng:weighted-clock}.  All estimates are uniform as $X_0$ tends
to zero.  Killed paths contribute nothing to either integral.
\end{proof}

\subsection{Shell summation and integration over physical area}

\begin{proof}[Proof of Proposition~\ref{eng:energy}]
Fix $t>0$, and divide $\Hh$ into the dyadic shells
\[
 A_j(t)=\{v\in\Hh:2^j\sqrt t\le|v|<2^{j+1}\sqrt t\},
 \qquad j\in\mathbb Z.
\]
Write $\vartheta=\arg v$ and
$L_j(t)=\int_{A_j(t)}\mu_t(v)|v|^{-2}\,\dd A(v)$.
Cauchy--Schwarz with the two factors
$\mu_t\sin^p\vartheta/|v|$ and
$\sin^{-p}\vartheta/|v|$ gives
\begin{align}
 L_j(t)^2
 &\le
 \left(\int_{A_j(t)}
       \frac{\mu_t(v)^2\sin^{2p}\vartheta}{|v|^2}\,\dd A(v)\right)
 \left(\int_{A_j(t)}
       \frac{\sin^{-2p}\vartheta}{|v|^2}\,\dd A(v)\right)\notag\\
 &\le K_p\int_{A_j(t)}
       \frac{\mu_t(v)^2\sin^{2p}\vartheta}{|v|^2}\,\dd A(v).
 \label{eng:shell-cs}
\end{align}
Indeed the second factor is exactly
\[
 (\log2)\int_0^\pi(\sin\vartheta)^{-2p}\,\dd\vartheta<\infty,
\]
because $2p<1$.  This is where the small angular power is retained.

For the $\epsilon$ chosen in \eqref{eng:epsilon-choice}, weighted
Cauchy--Schwarz over a finite set of shell indices gives
\[
 \left(\sum_j L_j(t)\right)^2
 \le\left(\sum_{j\in\mathbb Z}2^{-\epsilon|j|}\right)
      \sum_j2^{\epsilon|j|}L_j(t)^2.
\]
On $A_j(t)$, the weight $2^{\epsilon|j|}$ is bounded by a constant
times
$ (|v|/\sqrt t)^\epsilon+(\sqrt t/|v|)^\epsilon$.
Increasing the finite sets of indices and using monotone convergence
therefore gives the extended-integral inequality
\begin{equation}\label{eng:shell-global}
 J_t^2\le K\int_{\Hh}
   \frac{\mu_t(v)^2\sin^{2p}(\arg v)}{|v|^2}
   \left[
     \left(\frac{|v|}{\sqrt t}\right)^\epsilon
    +\left(\frac{\sqrt t}{|v|}\right)^\epsilon
   \right]\dd A(v).
\end{equation}

Under $w=\psi_t(v)$ the area formula is
\[
 \mu_t(v)^2\,\dd A(v)=m^2\,\dd A(w).
\]
The angular coordinate and radial coordinate in
\eqref{eng:shell-global} become $X_t(w)$ and $r_t(w)$ from
\eqref{eng:point-variables}.  Thus Tonelli's theorem gives
\begin{align}
 \E\int_0^\infty J_t^2\,\dd t
 &\le K m^2\int_\Omega
   \E\int_0^{T_w}
      X_t(w)^{2p}
      \left[
       \left(\frac{r_t(w)}{\sqrt t}\right)^\epsilon
       +\left(\frac{\sqrt t}{r_t(w)}\right)^\epsilon
      \right]\frac{\dd t}{r_t(w)^2}\,\dd A(w)\notag\\
 &\le K m^2\area(\Omega),
 \label{eng:area-conclusion}
\end{align}
where the last inequality is precisely
Lemma~\ref{eng:weighted} for each fixed physical point, followed by its
point-clock change of variables.  The constants are independent of
$w$.  There is no need to choose a common exceptional event for all
physical points: Tonelli applies to the nonnegative jointly measurable
integrands and the uniform expected bound for each fixed point.  The
single time $t=0$ has no effect on the integral.  This proves
\eqref{eng:energy-bound} and the almost-everywhere finiteness assertion.
\end{proof}

\begin{remark}\label{eng:scope}
The estimate holds on the entire capacity half-line.  Neither an
exponential-martingale estimate nor any absolute continuity with a
Brownian-driven law was used in obtaining it.  Its geometric inputs
remain uniform after live restarts, and every conformal-radius infimum
was taken before swallowing.  These two features are what make the
estimate available for the subsequent potential calculus.
\end{remark}
